\documentclass{article}
\usepackage{siunitx}
\usepackage{booktabs}
\usepackage{graphicx}
\usepackage{float}
\usepackage[version=4]{mhchem}
\usepackage{gensymb}
\usepackage{textcomp}
\sisetup{separate-uncertainty=true, per-mode=symbol}

\title{Lab Report: Specific Heat of Metals}
\author{Jane Doe \and John Roe}

\begin{document}
\maketitle

\section{Introduction}
We measured the specific heat of aluminium at \SI{25}{\celsius} and \SI{9.81}{\metre\per\second\squared}.
The sample mass was \qty{125.3(2)}{\gram} and the water volume \qty{250}{\milli\litre}.
The thermometer resolution is \num{0.1}; Avogadro's number is \num{6.022e23}.
Angles were measured at \ang{30} and \ang{12;30;0}. Temperatures rose \SIrange{20}{80}{\celsius}.
Room temperature was 22\degree C, or 22\,\textdegree C, about \SI{295}{\kelvin}.
The unit of power is \si{\watt} and of energy \unit{\kilo\joule\per\mole}.

The reaction is \ce{2H2 + O2 -> 2H2O} and the salt \ce{CuSO4.5H2O}.

\section{Results}
\begin{table}[H]
  \centering
  \caption{Measured specific heats.}
  \label{tab:heat}
  \begin{tabular}{l S[table-format=3.1] S[table-format=1.3]}
    \toprule
    {Metal} & {Mass (\si{\gram})} & {$c$ (\si{\joule\per\gram\per\kelvin})} \\
    \midrule
    Aluminium & 125.3 & 0.897 \\
    Copper    & 210.0 & 0.385 \\
    Iron      & 180.5 & 0.449 \\
    \bottomrule
  \end{tabular}
\end{table}

Using $Q = mc\Delta T$ with $\Delta T = \SI{12.5}{\kelvin}$, we find $Q = \SI{1.40e3}{\joule}$.
Table~\ref{tab:heat} shows values within \SI{2}{\percent} of literature, \tablenum{0.897}.

\end{document}
